\section{Sistem Persamaan, Prosedur Aljabar}

\begin{outcome}
\begin{enumerate}
\item[A.] Gunakan operasi elementer untuk menemukan solusi suatu sistem persamaan linear.

\item[B.] Temukan \ef\;dan \rref\;suatu matriks.

\item[C.] Tentukan dari \ef{} suatu sistem persamaan linear apakah sistem tersebut tidak memiliki solusi, memiliki
solusi unik, atau memiliki tak berhingga banyak solusi.

\item[D.] Selesaikan suatu sistem persamaan menggunakan eliminasi Gauss dan eliminasi Gauss--Jordan.

\item[E.] Modelkan suatu sistem fisik dengan persamaan linear, lalu selesaikan.
\end{enumerate}

\end{outcome}

Kita telah membahas secara mendalam representasi grafis sistem persamaan serta cara menemukan solusi yang mungkin
secara grafis. Sekarang perhatian kita beralih ke pengolahan sistem secara aljabar.

\begin{definition}{Sistem Persamaan Linear}\label{def:systemoflinearequations}
Suatu \textbf{sistem persamaan linear} \index{sistem persamaan} adalah daftar persamaan,
\begin{equation*}
\begin{array}{c}
a_{11}x_{1}+a_{12}x_{2}+\cdots +a_{1n}x_{n}=b_{1} \\
a_{21}x_{1}+a_{22}x_{2}+\cdots +a_{2n}x_{n}=b_{2} \\
\vdots \\
a_{m1}x_{1}+a_{m2}x_{2}+\cdots +a_{mn}x_{n}=b_{m}
\end{array}
\end{equation*}
dengan $a_{ij}$ dan $b_{j}$ berupa bilangan real. Bentuk di atas merupakan sistem
yang terdiri atas $m$ persamaan dalam $n$ variabel, $x_{1},x_{2}\cdots ,x_{n}$.
 Secara lebih ringkas, dengan menggunakan
notasi penjumlahan, bentuk di atas dapat ditulis sebagai
\begin{equation*}
\sum_{j=1}^{n}a_{ij}x_{j}=b_{i},
\text{ }i=1,2,3,\cdots ,m
\end{equation*}
Terakhir, kita dapat menuliskannya dalam notasi matriks sebagai
$$
A \x = \b 
$$
dengan
$$
A = \begin{bmatrix}
a_{11} & a_{12} & \dots & a_{1n}\\
a_{21} & a_{22} & \dots & a_{2n}\\
\vdots\\
a_{m1} & a_{m2} & \dots & a_{mn}
\end{bmatrix}, \ \ \ \x = \begin{bmatrix} x_1 \\ x_2 \\ \vdots \\ x_n\end{bmatrix}, \ \ \text{ dan } \ \ \b = \begin{bmatrix} b_1 \\ b_2 \\ \vdots \\ b_m\end{bmatrix}.
$$
\end{definition}

Ukuran relatif $m$ dan $n$ tidak penting di sini. Perhatikan bahwa kita
memperbolehkan $a_{ij}$ dan $b_{j}$ berupa sembarang bilangan real. Kita juga dapat
menyebut bilangan-bilangan ini \textbf{skalar} \index{skalar}. Istilah ini akan digunakan
di seluruh teks, jadi ingatlah bahwa istilah
\textbf{skalar} hanya berarti bahwa kita bekerja dengan bilangan real.

Sekarang, misalkan kita memiliki suatu sistem dengan $b_{i} = 0$ untuk setiap $i$. Dengan kata lain, setiap
persamaan sama dengan $0$. Ini merupakan jenis sistem khusus.

\begin{definition}{Sistem Persamaan Homogen}\label{def:homogeneoussystem}
Suatu sistem persamaan disebut \textbf{homogen} \index{sistem persamaan!homogen} jika setiap persamaan dalam
sistem tersebut sama dengan $0$. Sistem homogen memiliki bentuk
\begin{equation*}
\begin{array}{c}
a_{11}x_{1}+a_{12}x_{2}+\cdots +a_{1n}x_{n}= 0 \\
a_{21}x_{1}+a_{22}x_{2}+\cdots +a_{2n}x_{n}= 0  \\
\vdots \\
a_{m1}x_{1}+a_{m2}x_{2}+\cdots +a_{mn}x_{n}= 0 
\end{array}
\end{equation*}
dengan $a_{ij}$ berupa skalar dan $x_{i}$ berupa variabel.
\end{definition}

Ingat kembali dari bagian sebelumnya bahwa tujuan kita ketika bekerja dengan sistem persamaan linear
adalah menemukan titik perpotongan persamaan-persamaan tersebut ketika digambarkan. Dengan kata lain, kita
mencari solusi sistem. Sekarang kita ingin menemukan solusi-solusi tersebut secara aljabar. Kita ingin menemukan nilai
$ x_{1},\cdots ,x_{n} $ yang memenuhi semua
persamaan. Jika himpunan nilai semacam itu ada, kita menyebut $\left( x_{1},\cdots ,x_{n}\right)$ sebagai \textbf{himpunan solusi}. \index{sistem persamaan!himpunan solusi}

Ingat kembali pembahasan di atas mengenai jenis-jenis solusi yang mungkin.
Kita akan melihat bahwa sistem persamaan linear memiliki satu solusi unik, tak berhingga banyak solusi,
atau tidak memiliki solusi. Perhatikan definisi berikut.

\begin{definition}{Sistem Konsisten dan Tak Konsisten}\label{def:consistentandinconsistent}
Suatu sistem persamaan linear disebut
\index{sistem konsisten} \textbf{konsisten} jika terdapat sedikitnya satu solusi. Sistem tersebut
disebut
\index{sistem tak konsisten} \textbf{tak konsisten} jika tidak memiliki solusi.
\end{definition}

Jika setiap persamaan dianggap sebagai syarat yang harus
dipenuhi oleh variabel, konsisten berarti terdapat suatu pilihan nilai
variabel yang dapat memenuhi \textbf{semua} syarat. Tak konsisten berarti
tidak ada pilihan nilai variabel yang dapat memenuhi semua syarat.

Bagian-bagian berikut memberikan metode untuk menentukan apakah suatu sistem konsisten atau tak konsisten, serta untuk menemukan
solusinya jika solusi tersebut ada.
